\section{Merge Intervals}
\label{sec:merge-intervals}

\problemstatement{We are given an array of intervals
$\textit{intervals}[i] = [\textit{start}_i, \textit{end}_i]$, in no
particular order, and possibly overlapping with one another. We must merge
all overlapping intervals and return an array of the resulting
non-overlapping intervals that together cover exactly the same set of
points as the original collection.}

\begin{patternbox}
\emph{``Merge all overlapping intervals''} with \textbf{no} guarantee that
the input is already sorted is the unsorted counterpart of
Section~\ref{sec:insert-interval}. The keyword shift -- the absence of a
sortedness guarantee -- is the signal that we must add a sorting step
ourselves before the same kind of single linear merge pass can apply.
Whenever you see ``merge/combine all overlapping intervals'' with arbitrary
input order, the standard recipe is always: \emph{sort by start time, then
sweep once, extending or closing off a current ``running'' merged
interval}.
\end{patternbox}

\subsection*{Understanding the problem carefully}

Once the intervals are sorted by start time, two consecutive intervals (in
this sorted order) either overlap or they do not. Crucially, after sorting
by start time, if interval $j$ does not overlap the currently-growing merged
interval, no interval appearing even later (with an even larger start time)
can overlap it either, \emph{unless} merging interval $j$ itself first
extended the running interval far enough to reach that later interval --
but if $j$ did not overlap and was therefore not merged, the running
interval's end did not change, so this concern does not arise. This
guarantees that a single forward sweep, maintaining only the one
currently-open merged interval, is sufficient; we never need to revisit an
interval once we have moved past it.

\begin{ideabox}[Greedy Choice]
Sort all intervals by start time. Initialize the answer with the first
interval as the current ``running'' interval. For each subsequent interval,
in sorted order, check whether its start is at most the end of the running
interval. If so, it overlaps: extend the running interval's end to the
maximum of the two ends. If not, the running interval is finished growing:
push it into the answer, and start a new running interval from the current
one.
\end{ideabox}

\subsection*{Why sorting first makes the greedy sweep correct}

This is, in essence, the merging phase of Section~\ref{sec:insert-interval}
applied repeatedly, once for every interval, rather than just once for a
single new interval. The correctness argument is the same: because we
process intervals in increasing order of start time, the moment an interval
fails to overlap the current running interval, every future interval also
has a start time at least as large, hence also cannot reach backward to
overlap the running interval -- so the running interval can be safely
finalized and emitted. This sequential, never-look-back property is exactly
what makes a single $O(n)$ sweep (after the $O(n \log n)$ sort) sufficient,
with no need for any kind of nested loop checking all pairs of intervals
for overlap.

\subsection*{Algorithm}

\begin{enumerate}
  \item Sort $\textit{intervals}$ by start time, increasing.
  \item Initialize the result list with the first interval as the current
        running interval.
  \item For each subsequent interval $[\textit{s}, \textit{e}]$ in sorted
        order:
  \begin{enumerate}
    \item If $s \le \textit{lastEnd}$ (the end of the last interval placed
          in the result): extend that last interval's end to
          $\max(\textit{lastEnd}, e)$.
    \item Otherwise: append $[s, e]$ as a new entry in the result.
  \end{enumerate}
  \item Return the result.
\end{enumerate}

\begin{algorithm}[H]
\caption{Merge Intervals}
\begin{algorithmic}[1]
\Procedure{MergeIntervals}{\textit{intervals}}
  \State sort \textit{intervals} by start time, increasing
  \State $\textit{result} \gets [\textit{intervals}[0]]$
  \For{$i \gets 1$ to $n-1$}
    \State $\textit{last} \gets \textit{result}[\text{length}(\textit{result})-1]$
    \If{$\textit{intervals}[i].\textit{start} \le \textit{last}.\textit{end}$}
      \State $\textit{last}.\textit{end} \gets \max(\textit{last}.\textit{end}, \textit{intervals}[i].\textit{end})$
    \Else
      \State append $\textit{intervals}[i]$ to \textit{result}
    \EndIf
  \EndFor
  \State \Return \textit{result}
\EndProcedure
\end{algorithmic}
\end{algorithm}

\subsection*{Dry run}

\dryrun{Take $\textit{intervals} = [[1,3],[2,6],[8,10],[15,18]]$.}

Sorted by start time, the order is already
$[1,3],[2,6],[8,10],[15,18]$.

\begin{itemize}
  \item Initialize result $= [[1,3]]$.
  \item Process $[2,6]$: last in result is $[1,3]$, and $2 \le 3$, so
        overlap. Extend: $[1, \max(3,6)] = [1,6]$. Result $= [[1,6]]$.
  \item Process $[8,10]$: last in result is $[1,6]$, and $8 \le 6$? No.
        Append. Result $= [[1,6],[8,10]]$.
  \item Process $[15,18]$: last in result is $[8,10]$, and $15 \le 10$?
        No. Append. Result $= [[1,6],[8,10],[15,18]]$.
\end{itemize}

The final merged set is $[[1,6],[8,10],[15,18]]$.

\subsection*{C++ implementation}

\begin{lstlisting}
#include <bits/stdc++.h>
using namespace std;

vector<vector<int>> mergeIntervals(vector<vector<int>>& intervals) {
    if (intervals.empty()) return {};

    sort(intervals.begin(), intervals.end());

    vector<vector<int>> result;
    result.push_back(intervals[0]);

    for (int i = 1; i < (int)intervals.size(); i++) {
        if (intervals[i][0] <= result.back()[1]) {
            result.back()[1] = max(result.back()[1], intervals[i][1]);
        } else {
            result.push_back(intervals[i]);
        }
    }
    return result;
}

int main() {
    vector<vector<int>> intervals = {{1,3},{2,6},{8,10},{15,18}};
    auto merged = mergeIntervals(intervals);

    for (auto& iv : merged) {
        cout << "[" << iv[0] << "," << iv[1] << "] ";
    }
    cout << endl;
    return 0;
}
\end{lstlisting}

\begin{complexitybox}
\textbf{Time Complexity.} Sorting the $n$ intervals by start time costs
$O(n \log n)$. The subsequent single sweep costs $O(n)$. The overall time
complexity is
\[
  O(n \log n).
\]
\textbf{Space Complexity.} The result list holds at most $n$ intervals, so
the space complexity is $O(n)$. If in-place sorting is used and we are
allowed to modify the input, no further auxiliary space beyond the output
is required.
\end{complexitybox}

\begin{takeawaybox}
``Merge overlapping intervals'' and ``insert one interval into an
already-sorted, non-overlapping set'' (Section~\ref{sec:insert-interval})
are really the same merging logic at two different scales. Once you notice
that a single new interval inserted into a sorted list is structurally
identical to a sorted list being built up one interval at a time, both
problems collapse into the same primitive: maintain one ``currently open''
interval and only close it off when the next candidate's start exceeds its
end.
\end{takeawaybox}
